
\section{\textbf{\texttt{get\_metadata()}} ORM Method}

\begin{frame}[fragile]{\textbf{\texttt{get\_metadata()}} Method}
	\begin{itemize}
		\item The \textbf{\texttt{get\_metadata()}} method returns technical metadata about records.
		\item It includes information such as XML IDs, create/write audit fields, and whether the record is a noupdate data record.
		\item It is mostly used by technical tools and the UI debug / metadata dialogs.
\begin{block}{Method Syntax}
\begin{lstlisting}
def get_metadata(self):
	return super().get_metadata()
\end{lstlisting}
	\begin{itemize}
		\item Important parts: \textbf{\texttt{get\_metadata}}, \textbf{\texttt{self}}.
	\end{itemize}
\end{block}
	\end{itemize}
\end{frame}

\begin{frame}[fragile]{Breaking Down the Method Signature}
	\begin{block}{}
		\begin{itemize}
			\item \textbf{\texttt{self}}: The recordset to inspect.
			\item No values dictionary is required.
			\item Works on one or many records; returns one metadata dict per record.
\begin{lstlisting}
student = self.env['student.student'].browse(15)
meta = student.get_metadata()
\end{lstlisting}
		\end{itemize}
	\end{block}
\end{frame}

\begin{frame}[fragile]{What Does \texttt{get\_metadata()} Return?}
	\begin{block}{}
		\begin{itemize}
			\item It returns a \textbf{list of dictionaries}.
\begin{lstlisting}
meta = self.env['student.student'].browse(15).get_metadata()
print(meta)
\end{lstlisting}
			Output is (example):
\begin{lstlisting}
[{
	'id': 15,
	'xmlid': False,
	'noupdate': False,
	'create_uid': (1, 'Administrator'),
	'create_date': '2026-07-01 08:12:33',
	'write_uid': (2, 'Ahmed Muhumed'),
	'write_date': '2026-07-20 14:05:10',
}]
\end{lstlisting}
		\end{itemize}
	\end{block}
\end{frame}

\begin{frame}[fragile]{Important Metadata Keys}
	\begin{itemize}
		\item \texttt{id} -- database ID
		\item \texttt{xmlid} -- external ID if the record came from XML/data files
		\item \texttt{noupdate} -- whether upgrades should skip updating this data record
		\item \texttt{create\_uid} / \texttt{create\_date} -- who created it and when
		\item \texttt{write\_uid} / \texttt{write\_date} -- who last modified it and when
	\end{itemize}
\end{frame}

\begin{frame}[fragile]{When to Use \texttt{get\_metadata()}}
	\begin{itemize}
		\item Debugging module data and external IDs.
		\item Checking whether a record is system/data XML versus user-created.
		\item Inspecting create/write audit information quickly.
\begin{lstlisting}
for row in self.env['student.student'].browse([15, 16]).get_metadata():
	print(row['id'], row['xmlid'], row['write_date'])
\end{lstlisting}
	\end{itemize}
\end{frame}

\begin{frame}[fragile]{Method Call}
\begin{block}{Method Call}
\begin{lstlisting}
metadata = self.env['student.student'].browse(15).get_metadata()
\end{lstlisting}
	\begin{itemize}
		\item Here, \textbf{\texttt{.get\_metadata()}} is the \textbf{method call}.
		\item Return type reminder: \textbf{list of dicts}.
	\end{itemize}
\end{block}
\end{frame}
